\subsection{Classical Baselines on the Vesta IEEE--CIS Dataset}
\label{sec:vesta-classical-baselines}

We evaluate class-weighted logistic regression, random forest, XGBoost, and LightGBM on the IEEE--CIS fraud-detection data supplied by Vesta Corporation. The public training mirror contains 590,540 labeled e-commerce transactions and 20,663 fraud cases (3.499\%). Transaction and identity attributes are already merged. After removing 23 entirely missing columns, 432 predictors remain. To make the four-model and dimensionality-reduction experiment locally reproducible, we draw a deterministic stratified sample of 100,000 transactions, preserving 3,499 fraud cases. A stratified 80/20 holdout with seed 42 produces 80,000 training observations and 20,000 test observations, including 700 test frauds.

Numerical features are median-imputed; categorical features are most-frequent-imputed and ordinal-encoded. All preprocessing is fitted on the training partition only. Models use class weighting and a fixed 0.5 threshold. We evaluate both all 432 retained features and an eight-component PCA representation for a circuit-compatible comparison. PR-AUC is the primary metric because accuracy is inflated by the 96.501\% non-fraud majority.

\begin{table}[ht]
\centering
\caption{Classical holdout results on a stratified 100,000-row Vesta sample.}
\label{tab:vesta-classical-results}
\resizebox{\textwidth}{!}{%
\begin{tabular}{llrrrrrrrr}
\hline
Input & Model & Accuracy & Bal. Acc. & Precision & Recall & F1 & MCC & ROC-AUC & PR-AUC \\
\hline
All 432 & Logistic regression & 0.8345 & 0.7704 & 0.1367 & 0.7014 & 0.2288 & 0.2589 & 0.8462 & 0.3941 \\
All 432 & Random forest & \textbf{0.9744} & 0.6660 & \textbf{0.8357} & 0.3343 & \textbf{0.4776} & \textbf{0.5192} & 0.8987 & 0.5502 \\
All 432 & XGBoost & 0.8981 & \textbf{0.8336} & 0.2221 & \textbf{0.7643} & 0.3442 & 0.3767 & 0.9112 & 0.5651 \\
All 432 & LightGBM & 0.9200 & 0.8326 & 0.2672 & 0.7386 & 0.3924 & 0.4135 & \textbf{0.9154} & \textbf{0.5994} \\
\hline
PCA-8 & Logistic regression & 0.8329 & 0.7055 & 0.1157 & 0.5686 & 0.1923 & 0.2002 & 0.7886 & 0.2078 \\
PCA-8 & Random forest & \textbf{0.9710} & 0.6229 & \textbf{0.7632} & 0.2486 & \textbf{0.3750} & \textbf{0.4255} & \textbf{0.8322} & \textbf{0.4084} \\
PCA-8 & XGBoost & 0.8100 & 0.7419 & 0.1160 & \textbf{0.6686} & 0.1976 & 0.2215 & 0.8263 & 0.3555 \\
PCA-8 & LightGBM & 0.8440 & \textbf{0.7464} & 0.1353 & 0.6414 & 0.2234 & 0.2434 & 0.8288 & 0.3628 \\
\hline
\end{tabular}%
}
\end{table}

LightGBM provides the strongest ranking performance, attaining 59.94\% PR-AUC and 91.54\% ROC-AUC. XGBoost obtains the highest recall (76.43\%), detecting 535 of 700 test frauds, while random forest obtains the highest precision (83.57\%) but detects only 33.43\% of frauds. Reducing the input to eight principal components lowers the best PR-AUC from 59.94\% to 40.84\%. Consequently, an eight-input quantum model must be compared against the PCA-8 classical models evaluated on the identical sample and split. These values are local holdout measurements, not Kaggle leaderboard scores; final evaluation should use forward-in-time validation based on \texttt{TransactionDT}.

